\section{Profilo OWL 2 ed Espressivit\`a DL}
\label{sec:profilo_dl}

Questa sezione colloca l'ontologia rispetto alla teoria delle Description Logic
e ai profili di OWL 2, e spiega perch\'e lo schema \`e progettato per OWL 2 DL
mentre il reasoning pratico usa il profilo OWL 2 RL (Sezione~\ref{sec:reasoning}).

\subsection{Espressivit\`a: $\mathcal{SRIN}^{(\mathcal{D})}$}

OWL 2 DL corrisponde alla logica descrittiva $\mathcal{SROIQ}^{(\mathcal{D})}$.
I costrutti effettivamente usati nello schema collocano l'ontologia nel frammento
$\mathcal{SRIN}^{(\mathcal{D})}$. La Tabella~\ref{tab:espressivita} associa a ogni
lettera il costrutto e il punto dello schema che lo richiede.

\begin{table}[h]
\centering
\small
\caption{Espressivit\`a DL dell'ontologia}
\label{tab:espressivita}
\begin{tabularx}{\textwidth}{l >{\raggedright\arraybackslash}p{4.8cm} >{\raggedright\arraybackslash}X}
\toprule
\textbf{Lettera} & \textbf{Costrutto} & \textbf{Uso nello schema} \\
\midrule
$\mathcal{ALC}$ & $\sqcap$, $\exists$, $\lnot$ (e $\forall$, $\sqcup$ per dualit\`a) &
  classi definite \texttt{VideoGame} $\sqcap$ $\exists p.C$; complemento in \texttt{StandaloneGame} \\
$\mathcal{S}$   & $\mathcal{ALC}$ + ruoli transitivi & \texttt{sequelOf}, \texttt{partOfFranchise}, \texttt{subsidiaryOf} \\
$\mathcal{R}$   & inclusioni complesse di ruoli, asimmetria, irriflessivit\`a &
  3 property chain; 13 propriet\`a asimmetriche e 18 irriflessive \\
$\mathcal{I}$   & ruoli inversi & 12 \texttt{owl:inverseOf}; 3 propriet\`a simmetriche ($R \equiv R^-$) \\
$\mathcal{N}$   & restrizioni di cardinalit\`a non qualificate &
  ${\le}1\,\texttt{developedBy}$, ${\le}1\,\texttt{availableOn}$; \texttt{reviewedGame} funzionale \\
$(\mathcal{D})$ & datatype, facet, data property funzionali &
  26 datatype property; \texttt{minInclusive} in \texttt{HighlyRatedGame} e \texttt{RecentGame} \\
\bottomrule
\end{tabularx}
\end{table}

Rispetto a $\mathcal{SROIQ}^{(\mathcal{D})}$ mancano due componenti:
\begin{itemize}[nosep]
  \item $\mathcal{O}$ (nominali): lo schema non usa \texttt{owl:oneOf} n\'e
    \texttt{owl:hasValue}; le categorie (generi, piattaforme, modalit\`a) sono modellate
    come classi e non come insiemi chiusi di individui;
  \item $\mathcal{Q}$ (cardinalit\`a qualificate): le due restrizioni di cardinalit\`a
    non specificano \texttt{owl:onClass}.
\end{itemize}
Gli assiomi \texttt{owl:hasKey} non modificano l'espressivit\`a: sono \emph{DL-safe},
cio\`e si applicano solo agli individui nominati (Sezione~\ref{sec:haskey}).

\subsection{Profili OWL 2: perch\'e serve il reasoning a due livelli}

OWL 2 definisce tre profili trattabili (EL, QL, RL). L'ontologia non appartiene a
nessuno di essi:
\begin{itemize}[nosep]
  \item \textbf{OWL 2 EL} esclude ruoli inversi, complemento e cardinalit\`a;
  \item \textbf{OWL 2 QL} esclude, tra l'altro, le property chain e la transitivit\`a;
  \item \textbf{OWL 2 RL} ammette $\exists p.C$ solo come \emph{sottoclasse} e non ammette
    datatype restriction n\'e complemento in posizione di sottoclasse.
\end{itemize}

Il caso di OWL 2 RL \`e quello rilevante per il progetto. Un'equivalenza
$X \equiv \texttt{VideoGame} \sqcap \exists p.C$ contiene due inclusioni:
\begin{itemize}[nosep]
  \item $\texttt{VideoGame} \sqcap \exists p.C \sqsubseteq X$ (condizione sufficiente):
    ammessa in RL, \`e quella che serve per \emph{classificare} (regole
    \texttt{cls-int1} e \texttt{cls-svf1});
  \item $X \sqsubseteq \exists p.C$ (condizione necessaria): non ammessa in RL, perch\'e
    richiederebbe di introdurre individui anonimi.
\end{itemize}
Un reasoner RL come \texttt{owlrl} applica quindi correttamente la direzione utile
alla classificazione e ignora l'altra. Restano invece fuori dal profilo, in entrambe
le direzioni, le datatype restriction (\texttt{HighlyRatedGame}, \texttt{RecentGame})
e il complemento usato come condizione sufficiente (\texttt{StandaloneGame}): per
questi costrutti il progetto usa il secondo livello di reasoning basato su SPARQL
CONSTRUCT.

Le cardinalit\`a massime ${\le}1$ in posizione di superclasse, al contrario, sono
ammesse in RL: la regola \texttt{cls-maxc2} identifica (\texttt{owl:sameAs}) i due valori
di un individuo che ne avesse pi\`u di uno. Non producono per\`o classificazioni,
perch\'e \texttt{GameWithSingleDeveloper} ed \texttt{ExclusiveRelease} hanno sole
condizioni necessarie.

\subsection{Vincoli globali di OWL 2 DL}

Per garantire la decidibilit\`a, OWL 2 DL impone vincoli globali sull'uso delle
propriet\`a. Una propriet\`a \`e \emph{non semplice} se \`e transitiva, se \`e
super-propriet\`a di una catena, oppure se \`e inversa o super-propriet\`a di una
propriet\`a non semplice. Nello schema le propriet\`a non semplici sono sette:
\texttt{sequelOf}, \texttt{hasSequel}, \texttt{partOfFranchise}, \texttt{subsidiaryOf},
\texttt{sharedFranchiseWith}, \texttt{sharesDeveloperWith}, \texttt{sharesPublisherWith}.

I vincoli rispettati dallo schema sono:
\begin{enumerate}[nosep]
  \item \textbf{Propriet\`a semplici}: nessuna propriet\`a non semplice compare in
    assiomi di asimmetria, irriflessivit\`a, funzionalit\`a (diretta o inversa),
    cardinalit\`a o \texttt{hasKey}. Per questo le tre transitive e \texttt{hasSequel}
    non sono dichiarate asimmetriche n\'e irriflessive (Sezione~\ref{sec:tbox_proprieta});
  \item \textbf{Regolarit\`a delle catene}: le tre catene hanno la forma
    $R_1 \circ R_2 \sqsubseteq S$ con $S$ distinta da $R_1$ e $R_2$ e senza dipendenze
    cicliche tra catene, quindi la gerarchia di ruoli \`e regolare;
  \item \textbf{Nessuna InverseFunctionalProperty su datatype property}: OWL 2 DL non
    prevede l'inversa di una data property; l'identit\`a \`e espressa con \texttt{hasKey}.
\end{enumerate}

La verifica \`e stata eseguita con uno script basato su \texttt{rdflib}, che calcola
l'insieme delle propriet\`a non semplici e ne controlla l'uso negli assiomi elencati:
il risultato \`e di 0 violazioni.

\subsection{Limiti Noti}

Per l'uso con un reasoner DL completo (HermiT, Pellet) restano due aspetti:
\begin{itemize}[nosep]
  \item \textbf{Datatype \texttt{xsd:date}}. \`E il range di \texttt{releaseDate},
    \texttt{reviewDate} ed \texttt{eventDate} e la base della restrizione di
    \texttt{RecentGame}, ma non fa parte della \emph{datatype map} di OWL 2, che per gli
    istanti temporali prevede solo \texttt{xsd:dateTime} e \texttt{xsd:dateTimeStamp}.
    HermiT lo rifiuta come datatype non supportato; la soluzione \`e migrare schema e
    dati a \texttt{xsd:dateTime};
  \item \textbf{Entit\`a esterne non dichiarate}. Le entit\`a usate negli allineamenti
    (\texttt{wd:*}, \texttt{dcterms:*}) non hanno una dichiarazione di tipo esplicita;
    strumenti come l'OWL API le dichiarano implicitamente in fase di caricamento.
\end{itemize}
